\chapter{Type System Reference}

LLVM IR is strongly typed. Every value, instruction, and function has an explicit type.

\section{Integer Types}

\textbf{Fixed-width integers:}
\begin{lstlisting}[style=ts]
import { IntType } from "@euriklis/llvm-ir";

const i1 = LLVMCodegen.types.i1;      // Boolean (1 bit)
const i8 = LLVMCodegen.types.i8;      // Byte (8 bits)
const i16 = LLVMCodegen.types.i16;    // Short (16 bits)
const i32 = LLVMCodegen.types.i32;    // Int (32 bits)
const i64 = LLVMCodegen.types.i64;    // Long (64 bits)
const i128 = LLVMCodegen.types.i128;  // Quad (128 bits)
\end{lstlisting}

\textbf{Arbitrary-width integers:}
\begin{lstlisting}[style=ts]
const i7 = new IntType({ bits: 7 });    // 7-bit integer
const i24 = new IntType({ bits: 24 });  // 24-bit integer (RGB color)
const i256 = new IntType({ bits: 256 }); // 256-bit integer (crypto)
// Generates: i7, i24, i256
\end{lstlisting}

LLVM supports integers from \texttt{i1} to \texttt{i2\^{}23-1} (8 million bits).

\section{Floating-Point Types}

\begin{lstlisting}[style=ts]
import { FloatType } from "@euriklis/llvm-ir";

const half = new FloatType("half");      // f16: 16-bit half precision
const bfloat = new FloatType("bfloat");  // bf16: Brain float (ML)
const float = new FloatType("float");    // f32: Single precision (IEEE 754)
const double = new FloatType("double");  // f64: Double precision
const fp128 = new FloatType("fp128");    // Quad precision (128-bit)
\end{lstlisting}

\textbf{When to use each:}
\begin{itemize}[noitemsep]
  \item \texttt{half (f16)}: GPU compute, ML inference (low precision, high throughput)
  \item \texttt{bfloat (bf16)}: ML training (Google TPUs, NVIDIA Ampere+)
  \item \texttt{float (f32)}: Default for graphics and general compute
  \item \texttt{double (f64)}: Scientific computing, high-precision math
  \item \texttt{fp128}: Rare (extreme precision needs)
\end{itemize}

\section{Pointer Types}

In LLVM IR, pointers are opaque (no element type in modern LLVM):

\begin{lstlisting}[style=ts]
const ptr = LLVMCodegen.types.ptr;  // Opaque pointer
// Generates: ptr

// Address space variant
const ptr_addrspace3 = new PointerType(3);  // ptr addrspace(3)
// Generates: ptr addrspace(3) (CUDA shared memory)
\end{lstlisting}

\section{Array Types}

\begin{lstlisting}[style=ts]
import { ArrayType } from "@euriklis/llvm-ir";

const intArray = new ArrayType({
  elementType: i32,
  size: 10
});
// Generates: [10 x i32]

const floatMatrix = new ArrayType({
  elementType: new ArrayType({ elementType: float, size: 4 }),
  size: 4
});
// Generates: [4 x [4 x float]] (4x4 matrix)
\end{lstlisting}

\section{Struct Types}

\begin{lstlisting}[style=ts]
import { StructType } from "@euriklis/llvm-ir";

// Named struct
const Point = new StructType({
  name: "Point",
  fields: [
    { type: float },
    { type: float },
    { type: float }
  ]
});
// Generates: %Point = type { float, float, float }

// Packed struct (no padding)
const PackedPoint = new StructType({
  name: "PackedPoint",
  fields: [
    { type: i8 },
    { type: i32 }
  ],
  packed: true
});
// Generates: %PackedPoint = type <{ i8, i32 }>
\end{lstlisting}

\section{Vector Types}

SIMD vectors for parallel operations:

\begin{lstlisting}[style=ts]
import { VectorType } from "@euriklis/llvm-ir";

// <4 x float> - 4-element float vector (SSE, NEON)
const vec4 = new VectorType({
  elementType: float,
  count: 4
});

// <8 x i32> - 8-element integer vector (AVX)
const vec8i = new VectorType({
  elementType: i32,
  count: 8
});

// Vector addition
const vadd = new BinaryInst({
  name: "vsum",
  opcode: LLVMCodegen.opcodes.fadd,
  type: vec4,
  lhs: "a",
  rhs: "b"
});
// Generates: %vsum = fadd <4 x float> %a, %b
// This adds 4 floats in parallel!
\end{lstlisting}

\section{Function Types}

\begin{lstlisting}[style=ts]
import { FunctionType } from "@euriklis/llvm-ir";

// int add(int, int)
const addFnType = new FunctionType({
  returnType: i32,
  paramTypes: [i32, i32]
});
// Type: i32 (i32, i32)

// void print(char*)
const printFnType = new FunctionType({
  returnType: voidType,
  paramTypes: [ptr]
});
// Type: void (ptr)

// Variadic: int printf(char*, ...)
const printfFnType = new FunctionType({
  returnType: i32,
  paramTypes: [ptr],
  isVarArg: true
});
// Type: i32 (ptr, ...)
\end{lstlisting}

\section{Void Type}

Special type for functions that don't return a value:

\begin{lstlisting}[style=ts]
const { void: voidType } = LLVMCodegen.types;

const func = new LLVMFunction({
  name: "doSomething",
  returnType: voidType
});
// Generates: define void @doSomething() { ... }
\end{lstlisting}

\section{Type Compatibility}

LLVM is strictly typed. You cannot mix types without explicit conversion:

\begin{lstlisting}[style=ts]
// ERROR: Cannot add i32 and i64
const wrong = new BinaryInst({
  opcode: "add",
  type: i32,
  lhs: "x",  // i32
  rhs: "y"   // i64 - type mismatch!
});

// CORRECT: Convert first
const y64 = new ZExtInst({
  name: "y_extended",
  from: "y",
  fromType: i32,
  toType: i64
});
const correct = new BinaryInst({
  opcode: "add",
  type: i64,
  lhs: y64,
  rhs: "x"
});
\end{lstlisting}

